% MT013 / explanation — canonical editable content.
这个目标同时出现L1和L2，但它们作用的对象不同。预测值\(w^Tx_i=\sum_jw_jx_{ij}\)把各项输入乘上对应权重后相加。残差\(r_i=w^Tx_i-y_i\)就是预测值减真实值；\(|r_i|\)衡量预测错了多少，\(\|w\|_2^2=\sum_jw_j^2\)衡量权重整体有多大。\(\arg\min_w\)表示“找出使括号内总量最小的那组权重”，通常取\(\lambda\ge0\)。


\textbf{残差使用绝对损失，减轻极端输出误差的影响。} 误差从10变成20，绝对损失从10变20，平方损失却从100变400。平方损失会越来越强地要求模型照顾这个极端点；绝对损失在正负两侧的斜率大小都是1，随误差增加不继续放大。


这里“影响较小”不是“大误差与小误差罚得一样”。残差20仍比残差10有更大绝对损失，只是其增长速度不再随残差增加。对输出\(y\)异常通常更稳健；若输入\(x\)本身极大，参数梯度还会乘上\(x\)，所以不是对所有异常输入都免疫。


\textbf{权重使用L2，限制过大的系数。} \(\lambda\)越大，模型越不愿为降低训练误差而使用大权重，有助于稳定拟合。惩罚过强也可能欠拟合，强度仍需验证。


LASSO对\textbf{权重}取绝对值以鼓励精确零系数；本题绝对值在\textbf{残差}上，权重惩罚为平方，所以通常只有收缩作用，不自动得到稀疏特征选择。在\(\lambda\ge0\)时目标是凸的，即局部最小值也是全局最小值；但残差为0处有尖点，不能处处直接用普通导数。


\textbf{出处：} 2020B期中Quiz Q13，原卷第6页（10分）


\textbf{答案依据：} 2020B期中 Q13 随卷／配套参考答案；条件与解释经复核。
